\newpage
\phantomsection
\label{prf:falling-factorial-rule}
\LRAProofFor{thm:falling-factorial-rule}

\begin{remark*}[Return]
\hyperref[thm:falling-factorial-rule]{Return to Theorem}
\end{remark*}

\begin{theorem*}[Difference Rule for Falling Factorials]
For \(k \ge 1\),
\[
\Delta\, x^{\underline{k}} = k\,x^{\underline{k-1}}.
\]
\end{theorem*}

\begin{proof}[Professional Standard Proof]
\LRAProofBodyStart
Observe that
\[
(x+1)^{\underline{k}} = (x+1)\,x\,(x-1)\cdots(x-k+2).
\]
Therefore
\[
\Delta\, x^{\underline{k}}
= (x+1)^{\underline{k}} - x^{\underline{k}}
= x(x-1)\cdots(x-k+2)\bigl[(x+1)-(x-k+1)\bigr].
\]
The bracket simplifies to \(k\), and the remaining product is
\(x^{\underline{k-1}}\). Thus
\[
\Delta\, x^{\underline{k}} = k\,x^{\underline{k-1}}. \qedhere
\]
\end{proof}

\begin{proof}[Detailed Learning Proof]
\LRAProofBodyStart
Observe that
\[
(x+1)^{\underline{k}} = (x+1)\,x\,(x-1)\cdots(x-k+2).
\]
Therefore
\[
\Delta\, x^{\underline{k}}
= (x+1)^{\underline{k}} - x^{\underline{k}}
= x(x-1)\cdots(x-k+2)\bigl[(x+1)-(x-k+1)\bigr].
\]
The bracket simplifies to \(k\), and the remaining product is
\(x^{\underline{k-1}}\). Thus
\[
\Delta\, x^{\underline{k}} = k\,x^{\underline{k-1}}. \qedhere
\]
\end{proof}

\begin{remark*}[Proof structure]
Factor the two falling factorials by their shared descending product and
simplify the remaining linear difference.
\end{remark*}

\begin{dependencies}
\begin{itemize}
  \item \hyperref[def:falling-factorial]{Falling Factorial}
  \item \hyperref[def:forward-difference]{Forward Difference Operator}
\end{itemize}
\end{dependencies}
\clearpage
